% SPDX-FileCopyrightText: © 2026 Idiap Research Institute <contact@idiap.ch>
%
% SPDX-FileContributor: Andre Freitas <andre.freitas@idiap.ch>
% Neuro-symbolic AI Group
%
% SPDX-FileContributor: Danilo Gusicuma <danilo.gusicuma@idiap.ch>
% Neuro-symbolic AI Group
%
% SPDX-License-Identifier: MIT
% Source-bound conceptual exposition; aliases are local to each subsection.
\subsection{Cohere Fast: deciding whether a scientific target is usable}
\label{app:granular-cohere_fast}
\texttt{cohere\_fast} separates identifying a scientific target from deciding whether inquiry can responsibly begin. Its two occurrences do not yet explain or prove the target. The first, \texttt{pre:\_} in the original description, projects a domain target into a scope activation. The native name, \texttt{Unresolved\allowbreak Cohere\allowbreak Scope\allowbreak Projection}, preserves this missing mechanism rather than classifying it as abduction. Intake calibration then assesses that activation and returns a readiness disposition. The relation \texttt{Scope\allowbreak Activation\allowbreak Enables\allowbreak Cohere\allowbreak Intake} expresses a prerequisite: an identifiable activation enables assessment, but provides no evidence that the target claim is true.
\[
 p_{\_}:T\to\operatorname{Generated}[A],\qquad
 e:A\to\operatorname{Generated}[S],\qquad s=e(p_{\_}(t).\mathrm{value}).
\]
Here $T$ denotes \texttt{Cohere\allowbreak Domain\allowbreak Target}, $A$ denotes \texttt{Cohere\allowbreak Scope\allowbreak Activation}, and $S$ denotes \texttt{Cohere\allowbreak Intake\allowbreak Disposition}. The displayed \texttt{.value} means a successfully obtained generated payload; it is neither a proof nor a rule for ignoring execution failure. No multi-input join occurs.

\paragraph{Source-prescribed conceptual work.}
The Yang--Mills bindings instantiate a domain-specific programme rather than a generic default. Its target-freezing procedure separates theory existence, Hamiltonian spectral exclusion, clustering, confinement, and area law; these are different conclusions with different burdens. The representation procedures then make explicit which observable, state, dynamics, and limit structures a proposed bridge preserves. The central question becomes which implication in the proposed programme remains unsupported, not how many steps have been completed. The Cohere workflow contract supplies a child-stage transition structure, while individual skill bindings provide methods within those stages; the binding list alone establishes no execution order. In this reading, intake identifies the precise programme whose claims will be investigated, and later stages can return a conditional result or an obstruction without pretending to settle its ultimate target.

As a constructed Yang--Mills example, a proposal starts from decay of a finite set of Wilson-loop observables and announces a full mass gap. Intake identifies the intended pure-gauge target and the distinct missing claims: observable completeness, physical-unit uniformity, continuum existence, and reconstruction. The useful outcome is a programme of separate obligations instead of accepting the announced implication. No decay or mass-gap result is asserted by this example.

The original policy sends ready intake to \texttt{cohere\_fast\_scope\_and\_ground} and stops blocked or stalled work. The native reasoning declaration only returns $S$; it does not invoke that next workflow or construct its grounding request. Its module-level input adapter remains separate from the policy arrow. Thus the workflow's conceptual role is delimiting the question on which later reasoning can operate, with an explicit native realization gap. This gap does not imply that the source procedures lack conceptual content.

% SPDX-FileCopyrightText: © 2026 Idiap Research Institute <contact@idiap.ch>
%
% SPDX-FileContributor: Andre Freitas <andre.freitas@idiap.ch>
% Neuro-symbolic AI Group
%
% SPDX-FileContributor: Danilo Gusicuma <danilo.gusicuma@idiap.ch>
% Neuro-symbolic AI Group
%
% SPDX-License-Identifier: MIT
% Generated from step-explanations.json
\paragraph{Step-by-step reconstruction.}
A proposed Yang--Mills programme claims that decay of finitely many Wilson-loop observables establishes the full mass gap.
The following is a source-guided illustrative reconstruction. Internal procedures are prescriptions; the native declarations do not establish their execution.
\begin{enumerate}[leftmargin=*,label=\arabic*.]
\item \textbf{Identify the programme} (\texttt{scope\_\allowbreak{}activation}).
\begin{description}[style=standard,labelindent=0pt,leftmargin=5.5em,labelwidth=5em,labelsep=0.5em,align=left,font=\normalfont\bfseries,itemsep=0.15em,parsep=0pt,topsep=0.3em]
\item[Purpose] Separate the scientific target from its support.
\item[Input] Activated domain target; intended claim
\item[Method] The native scope producer is unresolved. Target-freezing instructions distinguish existence, spectral exclusion, clustering, confinement and area law before selecting a programme.
\item[Check] Different endpoints require different premises and cannot be merged by a shared label.
\item[Output] Intended activation: pure-gauge programme; separate target burdens
\item[Next] Intake calibration consumes the activation; its existence does not establish any target theorem.
\item[Example] Split the announcement into finite-observable decay, observable completeness, continuum existence, reconstruction and spectral exclusion instead of accepting one as all the others.
\end{description}
\item \textbf{Judge whether inquiry can begin} (\texttt{intake\_\allowbreak{}disposition}).
\begin{description}[style=standard,labelindent=0pt,leftmargin=5.5em,labelwidth=5em,labelsep=0.5em,align=left,font=\normalfont\bfseries,itemsep=0.15em,parsep=0pt,topsep=0.3em]
\item[Purpose] Assess whether the scoped question is usable.
\item[Input] Scope activation; domain; target
\item[Method] Assess whether the target and its boundaries are definite enough to investigate, retaining missing claims as obligations.
\item[Check] Readiness concerns a usable question, not support for its ultimate conclusion.
\item[Output] Intake disposition: ready, blocked or stalled
\item[Next] Original policy routes ready intake to grounding; the native return constructs no grounding request.
\item[Example] Proceed only with the separate obligations visible; an unspecified meaning of mass gap requires target clarification before scientific grounding.
\end{description}
\end{enumerate}
\paragraph{Changed-input contrast.}
If intake is instead asked only for a finite-regulator observable estimate, continuum reconstruction is not part of that narrower result. The activation must preserve the selected target rather than silently exchanging the two questions.


\subsection{Scope and grounding: restricting the admissible claim}
\label{app:granular-cohere_fast_scope_and_ground}
\texttt{cohere\_fast\_scope\_and\_ground} makes domain, target, sources, and maximum claim scope explicit before discovery proceeds. Its unresolved projection transforms a request containing domain, target, and source references into a grounding record that also names a claim ceiling. The subsequent evidential assessment asks whether that bounded grounding is usable. These are distinct responsibilities: constructing a scope description and judging its adequacy. The original ERFL description gives no reasoning operator for the former, retaining \texttt{pre:\_}; naming it a projection does not implement the source-interpretation procedure prescribed by its linked skills.
\[
 p_{\_}:Q\to\operatorname{Generated}[G],\qquad
 e:G\to\operatorname{Generated}[S],\qquad s=e(p_{\_}(q).\mathrm{value}).
\]
Here $Q$ denotes \texttt{Domain\allowbreak Scope\allowbreak Grounding\allowbreak Request}, $G$ denotes \texttt{Domain\allowbreak Scope\allowbreak Grounding}, and $S$ denotes \texttt{Scope\allowbreak Grounding\allowbreak Disposition}. The generated payload projection is conditional on successful acquisition. The single relation, \texttt{Domain\allowbreak Scope\allowbreak Grounding\allowbreak Enables\allowbreak Disposition}, licenses assessment of that grounding as a dependency; it is not an inference from source presence to claim truth. There is no join.

\paragraph{Source-prescribed conceptual work.}
In the Yang--Mills instantiation, target freezing first specifies the pure gauge theory, compact simple gauge group, dimension, axiomatic endpoint, vacuum sector, and regulator, volume, cutoff, and coupling quantifiers. It partitions premises from targets: assuming a positive physical correlation scale would conceal the very estimate being sought. Literature grounding binds each imported claim to a locatable primary source and retains its conventions and dependencies. The conceptual product is an admissible implication problem with a precise scope, not merely a source list. A strong-coupling regulated statement, an abstract transfer theorem, and a continuum Hamiltonian claim cannot be interchanged. These procedures prescribe what the native grounding record is meant to summarize; its string fields do not establish that the interpretation or preserved scope is correct.

In a constructed Yang--Mills case, a cited estimate holds at fixed lattice spacing in a strong-coupling regime. Grounding retains that regime and asks which constants depend on volume and cutoff. It does not relabel the estimate as a continuum mass gap. The exact open question is whether a suitable physical decay rate survives the required limits under a justified comparison, rather than whether the source contains the words ``mass gap''.

Original accepted status routes to adaptive frontier formation; blocked status stops. Native code returns only the disposition. The next workflow's identically named \texttt{Domain\allowbreak Scope\allowbreak Grounding} is module-local and has a different field set; neither matching names nor the policy transition supplies an adapter. A handoff must preserve the intended target and claim ceiling explicitly. The unresolved grounding operation and this cross-module boundary remain visible rather than being replaced by an invented executable chain.

% SPDX-FileCopyrightText: © 2026 Idiap Research Institute <contact@idiap.ch>
%
% SPDX-FileContributor: Andre Freitas <andre.freitas@idiap.ch>
% Neuro-symbolic AI Group
%
% SPDX-FileContributor: Danilo Gusicuma <danilo.gusicuma@idiap.ch>
% Neuro-symbolic AI Group
%
% SPDX-License-Identifier: MIT
% Generated from step-explanations.json
\paragraph{Step-by-step reconstruction.}
The source packet contains a strong-coupling estimate at fixed lattice spacing. The intended endpoint is a physical continuum spectral statement.
The following is a source-guided illustrative reconstruction. Internal procedures are prescriptions; the native declarations do not establish their execution.
\begin{enumerate}[leftmargin=*,label=\arabic*.]
\item \textbf{Construct admissible grounding} (\texttt{scope\_\allowbreak{}grounding}).
\begin{description}[style=standard,labelindent=0pt,leftmargin=5.5em,labelwidth=5em,labelsep=0.5em,align=left,font=\normalfont\bfseries,itemsep=0.15em,parsep=0pt,topsep=0.3em]
\item[Purpose] Ground a bounded implication problem in sources.
\item[Input] Grounding request: domain; target; source references
\item[Method] The native producer is unresolved. Linked instructions freeze theory, gauge group, quantifiers and endpoint, then bind imported claims to locatable sources with their conventions.
\item[Check] Do not assume the desired physical scale or erase a source theorem's regulator regime.
\item[Output] Intended grounding: fixed-spacing support; explicit claim ceiling
\item[Next] Disposition assessment receives the grounding; nominal fields do not ensure faithful source interpretation.
\item[Example] Retain the estimate's strong-coupling and fixed-spacing assumptions and mark dependence on volume and cutoff as questions to inspect.
\end{description}
\item \textbf{Assess grounding adequacy} (\texttt{grounding\_\allowbreak{}disposition}).
\begin{description}[style=standard,labelindent=0pt,leftmargin=5.5em,labelwidth=5em,labelsep=0.5em,align=left,font=\normalfont\bfseries,itemsep=0.15em,parsep=0pt,topsep=0.3em]
\item[Purpose] Assess source support for further inquiry.
\item[Input] Grounding record; target; claim ceiling
\item[Method] Compare imported assumptions with the requested endpoint and identify the missing implication between their regimes.
\item[Check] Fixed-regulator support must not be relabeled as a continuum Hamiltonian conclusion.
\item[Output] Grounding disposition: usable bounded support or specified blocker
\item[Next] Original policy routes accepted grounding onward; module-local type names still require a preserving adapter.
\item[Example] Accept the source as a restricted starting point while keeping physical-rate uniformity and continuum passage open; block a packet that silently treats them as proved.
\end{description}
\end{enumerate}
\paragraph{Changed-input contrast.}
A source theorem uniform in the required limits would change the grounding burden, but only after its assumptions and theory identity were checked. Matching terminology is never sufficient transfer evidence.


\subsection{Adaptive frontier: identifying the next bounded question}
\label{app:granular-cohere_fast_adaptive_frontier}
\texttt{cohere\_fast\_adaptive\_frontier} organizes what to investigate next within an existing claim boundary. The first occurrence projects a grounding into a frontier reference with the intended preservation of domain and claim ceiling. Its string fields do not enforce equality with the input across this projection. Its native name, \texttt{Unresolved\allowbreak Domain\allowbreak Frontier\allowbreak Projection}, reflects the original \texttt{pre:\_}: the declaration specifies the frontier interface, while the linked procedures below prescribe how alternatives and representation burdens are investigated. Evidential readiness then judges whether the proposed frontier is suitable for the derivation lane.
\[
 p_{\_}:G\to\operatorname{Generated}[F],\qquad
 e:F\to\operatorname{Generated}[S],\qquad s=e(p_{\_}(g).\mathrm{value}).
\]
Here $G$ denotes \texttt{Domain\allowbreak Scope\allowbreak Grounding}, $F$ denotes \texttt{Bounded\allowbreak Domain\allowbreak Frontier}, and $S$ denotes \texttt{Frontier\allowbreak Readiness\allowbreak Disposition}, all local nominal types. Every output remains generated; \texttt{.value} denotes successful payload access. \texttt{Domain\allowbreak Frontier\allowbreak Enables\allowbreak Readiness\allowbreak Assessment} records the prerequisite for the assessment, with no additional scientific warrant and no multi-input join.

\paragraph{Source-prescribed conceptual work.}
The linked Yang--Mills procedures construct alternatives along three connected axes. A regulator atlas compares carriers, measures, boundary conditions, gauge invariance, locality, reflection positivity, and closure under blocking; it chooses a primary representation and an adversarial alternative by their missing theorem obligations. Observable-core construction asks whether the selected observables separate states and generate the physical sectors needed by a spectral argument, then seeks a hidden-sector countermodel. The representation calculus specifies observable pullbacks, state pushforwards, their compatibility, and observable-indexed approximation errors. It distinguishes equivalence from coarse-graining or a conjectural correspondence. Restricted-model calibration tests these proposals on massive, massless, finite-volume, and hidden-sector cases and revises definitions that produce false positives. This is the substantive frontier-building work compressed by the native projection: discovering which representation and observable choices expose a tractable missing implication, while preserving what could invalidate transfer.

A constructed diagnostic example uses a direct sum of a gapped sector and an unobserved gapless sector. An observable family restricted to the first sector can show decay while failing to detect the second. The frontier therefore changes from ``fit a faster decay curve'' to ``establish density of the observable-generated physical subspace, or report a sectorial claim.'' This is a calibration countermodel prescribed by the source procedures, not a claim that the Yang--Mills theory has such a decomposition.

The original policy sends accepted frontiers to \texttt{cohere\_fast\_derivation\_and\_proof} and stops quarantined ones. Native reasoning returns the readiness disposition, not the frontier needed by that next input. Exposing or reconstructing the frontier for a subsequent call requires a separate handoff. The native declaration itself does not implement an adaptive search loop: ``adaptive'' names the regime's purpose, while realization of its prescribed conceptual work as this native occurrence remains unresolved.

% SPDX-FileCopyrightText: © 2026 Idiap Research Institute <contact@idiap.ch>
%
% SPDX-FileContributor: Andre Freitas <andre.freitas@idiap.ch>
% Neuro-symbolic AI Group
%
% SPDX-FileContributor: Danilo Gusicuma <danilo.gusicuma@idiap.ch>
% Neuro-symbolic AI Group
%
% SPDX-License-Identifier: MIT
% Generated from step-explanations.json
\paragraph{Step-by-step reconstruction.}
In an abstract spectral calibration model, a visible gapped sector yields decaying correlations while an unobserved gapless sector coexists with it. This tests an inference, not existence of a gapless Yang--Mills theory. The frontier must expose the missing completeness burden.
The following is a source-guided illustrative reconstruction. Internal procedures are prescriptions; the native declarations do not establish their execution.
\begin{enumerate}[leftmargin=*,label=\arabic*.]
\item \textbf{Build a discriminating frontier} (\texttt{bounded\_\allowbreak{}frontier}).
\begin{description}[style=standard,labelindent=0pt,leftmargin=5.5em,labelwidth=5em,labelsep=0.5em,align=left,font=\normalfont\bfseries,itemsep=0.15em,parsep=0pt,topsep=0.3em]
\item[Purpose] Expose the missing implication through representation questions.
\item[Input] Grounded domain; fixed claim ceiling
\item[Method] The native producer is unresolved. Linked methods compare regulators, observable coverage and representation maps, then use restricted massive, massless and hidden-sector cases to test false positives.
\item[Check] A representation earns a role by preserving required structure and exposing a falsifiable burden.
\item[Output] Intended frontier: observable-density obligation or sectorial conclusion
\item[Next] Readiness assesses this frontier; the signature does not implement adaptive search.
\item[Example] The hidden-sector construction changes the next question from fitting decay to proving that the chosen observable class sees every relevant physical sector.
\end{description}
\item \textbf{Judge frontier readiness} (\texttt{frontier\_\allowbreak{}readiness}).
\begin{description}[style=standard,labelindent=0pt,leftmargin=5.5em,labelwidth=5em,labelsep=0.5em,align=left,font=\normalfont\bfseries,itemsep=0.15em,parsep=0pt,topsep=0.3em]
\item[Purpose] Assess whether the frontier supports bounded derivation.
\item[Input] Frontier; preserved domain and claim ceiling
\item[Method] Check whether the chosen representation names the objects, required map and missing theorem that derivation must address.
\item[Check] A promising representation without its transfer obligation is not an unconditional route.
\item[Output] Frontier disposition: ready or quarantined
\item[Next] The native return is a disposition, not the frontier needed by the next workflow.
\item[Example] A precise density obligation can enter derivation as an open task; a claim that visible decay already proves global exclusion must be quarantined.
\end{description}
\end{enumerate}
\paragraph{Changed-input contrast.}
Once the observable-generated physical subspace is known to be dense and the extension estimates are justified, the next frontier can move to remaining transfer and limit obligations. Repeating the same visible-sector fit does not supply density.


\subsection{Derivation and proof: exposing a conditional result}
\label{app:granular-cohere_fast_derivation_and_proof}
\texttt{cohere\_fast\_derivation\_and\_proof} demands a derivation packet whose unresolved obligations remain attached. Despite the workflow name, its original first occurrence is \texttt{pre:\_}, not a declared deduction operator. The native \texttt{Unresolved\allowbreak Conditional\allowbreak Derivation\allowbreak Projection} specifies a transformation from a bounded frontier into a packet containing a derivation reference, open-obligations reference, domain, and claim ceiling. Evidential assessment then asks whether that packet is ready for audit; it returns no proof verdict.
\[
 p_{\_}:F\to\operatorname{Generated}[D],\qquad
 e:D\to\operatorname{Generated}[S],\qquad s=e(p_{\_}(f).\mathrm{value}).
\]
Here $F$ denotes \texttt{Bounded\allowbreak Domain\allowbreak Frontier}, $D$ denotes \texttt{Conditional\allowbreak Derivation\allowbreak Packet}, and $S$ denotes \texttt{Derivation\allowbreak Readiness\allowbreak Disposition}. The successful \texttt{.value} adapter does not certify the derivation. \texttt{Derivation\allowbreak Packet\allowbreak Enables\allowbreak Readiness\allowbreak Assessment} marks a dependency between packet and assessment, without converting readiness into entailment. There is no join of independent evidence here.

\paragraph{Source-prescribed conceptual work.}
The Yang--Mills proof bindings go substantially beyond packaging a conditional formula. Adaptive block domination fixes the exposure filtration and retained boundary information, bounds normalized probability on a good chart, and charges the complementary configurations through a multi-link estimate. Heat-kernel prefactors, chart mass, Haar Jacobians, overlap, and remainder oscillation enter the normalized comparison; a pointwise action estimate alone does not supply it. A state-weighted alternative proves a positive-kernel drift inequality $QV\leq q_VV$, controls the initial weight moment, and derives iteration before checking a strict cluster threshold. Retained-state sufficiency, positive-mass histories, and uniformity in volume and boundary conditions are independent obligations. Finite-block and enclosure candidates make parts of these burdens checkable; exact finite checks do not discharge full-measure, repeated-RG, continuum, or reconstruction hypotheses. The prescribed broader programme separates basin entry, controlled iteration, physical-unit decay, positivity, observable density, and continuum passage. These are implication dependencies to prove, not guaranteed successes or an execution order inferred from the skill list.

For a constructed finite diagnostic, the law $P(10)=P(01)=1/2$ has vanishing two-site intersection probability, yet $P(X_2=1\mid X_1=0)=1$. A small intersection moment therefore cannot replace the history-uniform conditional ceiling required by the domination argument. The next proof task is to identify an exposure state and prove its normalized conditional bound, not to declare domination from that moment. This source-prescribed counterexample tests a proposed proof mechanism; it proves no Yang--Mills, continuum, or Hamiltonian-gap conclusion.

Original accepted status routes to independent audit; quarantine stops. Native reasoning returns only readiness, whereas audit requires the actual conditional packet. That handoff is not executable merely because the workflow names are adjacent. Both the missing derivation producer and the preservation of open obligations must remain explicit when describing the larger chain.

% SPDX-FileCopyrightText: © 2026 Idiap Research Institute <contact@idiap.ch>
%
% SPDX-FileContributor: Andre Freitas <andre.freitas@idiap.ch>
% Neuro-symbolic AI Group
%
% SPDX-FileContributor: Danilo Gusicuma <danilo.gusicuma@idiap.ch>
% Neuro-symbolic AI Group
%
% SPDX-License-Identifier: MIT
% Generated from step-explanations.json
\paragraph{Step-by-step reconstruction.}
A proposed domination argument treats a small joint event probability as a history-uniform conditional ceiling. Test the finite law assigning probability one half each to outcomes 10 and 01.
The following is a source-guided illustrative reconstruction. Internal procedures are prescriptions; the native declarations do not establish their execution.
\begin{enumerate}[leftmargin=*,label=\arabic*.]
\item \textbf{Construct the conditional derivation} (\texttt{conditional\_\allowbreak{}derivation}).
\begin{description}[style=standard,labelindent=0pt,leftmargin=5.5em,labelwidth=5em,labelsep=0.5em,align=left,font=\normalfont\bfseries,itemsep=0.15em,parsep=0pt,topsep=0.3em]
\item[Purpose] Specify the mechanism’s required inequality and assumptions.
\item[Input] Bounded frontier; target; claim ceiling
\item[Method] The native producer is unresolved. Linked domination methods specify exposure history, retained boundary information and normalized conditional bounds; finite diagnostics attack the required implication.
\item[Check] Joint moments cannot replace history-uniform bounds; normalization and positive-probability conditioning must be justified.
\item[Output] Intended derivation packet: failed inference; open replacement obligation
\item[Next] Readiness assesses the packet; the declared projection is not a proof-producing algorithm.
\item[Example] For P(10)=P(01)=1/2, the two-site intersection probability is zero, but P(X2=1 given X1=0)=1. Thus the suggested conditional ceiling does not follow.
\end{description}
\item \textbf{Assess readiness for audit} (\texttt{derivation\_\allowbreak{}readiness}).
\begin{description}[style=standard,labelindent=0pt,leftmargin=5.5em,labelwidth=5em,labelsep=0.5em,align=left,font=\normalfont\bfseries,itemsep=0.15em,parsep=0pt,topsep=0.3em]
\item[Purpose] Assess whether the derivation is auditable.
\item[Input] Conditional packet; open-obligation reference
\item[Method] Check that the proposed inequality, retained state, assumptions and failure cases remain explicit rather than hidden by a certificate label.
\item[Check] A finite check cannot discharge full-measure, repeated-iteration or continuum hypotheses.
\item[Output] Derivation disposition: ready for scoped audit or quarantined
\item[Next] Original policy sends accepted packets to audit; native return omits the packet adapter.
\item[Example] Send the counterexample and the requirement for a normalized history-uniform bound to audit; withhold any domination conclusion inferred only from the zero intersection moment.
\end{description}
\end{enumerate}
\paragraph{Changed-input contrast.}
A proved normalized conditional bound for every admissible positive-probability exposure history would answer this particular objection. Small unconditional intersection moments alone do not, and even that bound would not settle continuum reconstruction.


\subsection{Independent audit: challenging a claim without losing it}
\label{app:granular-cohere_fast_independent_audit}
\texttt{cohere\_fast\_independent\_audit} has two declared reasoning mechanisms. Refutation examines a conditional derivation for a defeater and returns \texttt{Refuted}, \texttt{Not\allowbreak Refuted\allowbreak Within\allowbreak Scope}, or \texttt{Inconclusive}. Calibration then receives both the original derivation and that assessment. This join matters because the same adverse finding can refute an overstrong conclusion while leaving a narrower conditional implication intact. The relation \texttt{Calibrates\allowbreak Independent\allowbreak Audit\allowbreak Disposition} licenses use of the audit finding to assess the retained packet; it supplies no theorem independently of the finding.
\[
 r:D\to\operatorname{Generated}[V],\quad
 e:J\to\operatorname{Generated}[S].
\]
\[
 s=e\bigl(\{\mathrm{derivation}=d,\mathrm{refutation\_assessment}=r(d).\mathrm{value}\}\bigr).
\]
Here $D$ denotes \texttt{Conditional\allowbreak Derivation\allowbreak Packet}, $V$ denotes \texttt{Independent\allowbreak Refutation\allowbreak Assessment}, $J$ denotes \texttt{Independent\allowbreak Audit\allowbreak Disposition\allowbreak Input}, and $S$ denotes \texttt{Independent\allowbreak Audit\allowbreak Disposition}. Generated assessment access is a successful-result view. $J$ is a nominal record, not a native dependent pair enforcing the assessment's exact claim index.

\paragraph{Source-prescribed conceptual work.}
The linked Yang--Mills audit expands every cutoff, volume, distance, observable, and coupling quantifier, then seeks the first unproved implication in the complete proposed argument. It checks whether a bridge encodes its desired conclusion in its definition, whether constants remain uniform in physical units, and whether a representation change loses positivity, locality, or observable completeness. Limit interchanges require their own justification: a favorable order of limits is not automatically licensed. The audit distinguishes stochastic mixing, finite transfer spectra, sectorial decay, and a continuum Hamiltonian gap. Countermodels and the finite checkable core are used to attack particular steps rather than to cast a generic confidence vote. Its output is the strongest defensible statement together with the exact missing bridge and required downgrade language. This source-prescribed adversarial programme explains what the native refutation assessment summarizes; its name alone supplies no empirical independence guarantee.

In a constructed audit, a proposed spectral argument supplies temporal decay only for a finite observable family. The hidden-sector countermodel challenges the step from those correlations to the full spectrum. Refutation targets the missing density implication; calibration can retain a conditional or sectorial result while withholding the full claim. This differs from rejecting the decay calculation itself. Bounded non-refutation of the calculation cannot fill the missing density proof.

Original accepted status routes to manuscript release and quarantined status stops. Native code returns $S$ only. Release instead expects an \texttt{Independent\allowbreak Audit\allowbreak Packet} containing an audit reference and claim ceiling, so a separate construction must retain the reviewed substance. No unresolved \texttt{\_} occurrence appears inside this audit, but the cross-workflow packet adapter remains outside its signature.

% SPDX-FileCopyrightText: © 2026 Idiap Research Institute <contact@idiap.ch>
%
% SPDX-FileContributor: Andre Freitas <andre.freitas@idiap.ch>
% Neuro-symbolic AI Group
%
% SPDX-FileContributor: Danilo Gusicuma <danilo.gusicuma@idiap.ch>
% Neuro-symbolic AI Group
%
% SPDX-License-Identifier: MIT
% Generated from step-explanations.json
\paragraph{Step-by-step reconstruction.}
A conditional spectral argument relies on decay for finitely many observables and leaves density of their generated physical sector unproved.
The following is a source-guided illustrative reconstruction. Internal procedures are prescriptions; the native declarations do not establish their execution.
\begin{enumerate}[leftmargin=*,label=\arabic*.]
\item \textbf{Attack the complete implication} (\texttt{refutation\_\allowbreak{}assessment}).
\begin{description}[style=standard,labelindent=0pt,leftmargin=5.5em,labelwidth=5em,labelsep=0.5em,align=left,font=\normalfont\bfseries,itemsep=0.15em,parsep=0pt,topsep=0.3em]
\item[Purpose] Locate the calculation-to-endpoint inference gap.
\item[Input] Conditional derivation; quantifiers; open obligations
\item[Method] Expand volume, cutoff, distance and observable quantifiers; test hidden sectors, limit interchanges and definitions that encode the desired conclusion.
\item[Check] A defect must target a particular inference; bounded non-refutation is not full proof.
\item[Output] Refutation assessment: missing density implication or another scoped defeater
\item[Next] Calibration receives this assessment with the original derivation; the name independent proves no empirical independence.
\item[Example] As an abstract spectral countermodel, adjoin an orthogonal sector with positive energies approaching zero, invisible to all tested observables. Their decay survives, so the general global-gap inference fails; this does not construct a Yang--Mills counterexample.
\end{description}
\item \textbf{Retain the defensible claim} (\texttt{audit\_\allowbreak{}disposition}).
\begin{description}[style=standard,labelindent=0pt,leftmargin=5.5em,labelwidth=5em,labelsep=0.5em,align=left,font=\normalfont\bfseries,itemsep=0.15em,parsep=0pt,topsep=0.3em]
\item[Purpose] Calibrate the claim while preserving its valid part.
\item[Input] Original derivation + its refutation assessment
\item[Method] Compare the audited scope with the failure, identify the strongest retained statement and state the required downgrade.
\item[Check] Do not confuse refuting an overstrong extension with refuting its narrower conditional calculation.
\item[Output] Audit disposition: retained sectorial support; quarantined full-spectrum inference
\item[Next] Release requires a separate audit packet preserving this content; native disposition alone is insufficient.
\item[Example] Keep the visible-sector decay statement and make density an explicit missing bridge, preventing a favorable local audit from certifying a global mass gap.
\end{description}
\end{enumerate}
\paragraph{Changed-input contrast.}
If the packet states only visible-sector decay, the hidden-sector example does not refute that restricted calculation. It defeats the stronger global inference, so calibration must preserve the distinction.


\subsection{Manuscript release: preserving authority while presenting a result}
\label{app:granular-cohere_fast_manuscript_release}
\texttt{cohere\_fast\_manuscript\_release} asks whether a presentation stays within the authority of its source audit. Its first occurrence is intended to construct a presentation reference tied to the audit and its claim ceiling. String references alone do not enforce that cross-projection preservation. The original \texttt{pre:\_} and native \texttt{Unresolved\allowbreak Claim\allowbreak Bounded\allowbreak Presentation\allowbreak Projection} leave the native realization unresolved; the linked procedures prescribe the paper-construction work described below. The second occurrence calibrates release disposition from that presentation. It assesses an exposition, not a new scientific argument that strengthens its source claims.
\[
 p_{\_}:A\to\operatorname{Generated}[P],\qquad
 e:P\to\operatorname{Generated}[S],\qquad s=e(p_{\_}(a).\mathrm{value}).
\]
Here $A$ denotes \texttt{Independent\allowbreak Audit\allowbreak Packet}, $P$ denotes \texttt{Claim\allowbreak Bounded\allowbreak Presentation}, and $S$ denotes \texttt{Manuscript\allowbreak Release\allowbreak Disposition}. The payload adapter remains conditional on successful generation. The single relation \texttt{Claim\allowbreak Bounded\allowbreak Presentation\allowbreak Enables\allowbreak Release} expresses the prerequisite for the release assessment; it adds no evidence for the presented scientific proposition. There is no multi-input join: the source audit reference and ceiling must already be retained in $P$.

\paragraph{Source-prescribed conceptual work.}
The domain-private writing bindings divide presentation into organizing accepted results and explicit open obligations, assembling proof details without upgrading their status, and rendering the manuscript from the accepted evidence packets. Each procedure retains conventions, dependencies, attempted falsifiers, failed routes, and conditional branches. The resulting argument must show which theorem uses which hypothesis and where a finite certificate stops short of its intended application. This is more than checking that a file compiles: a reader must be able to distinguish established mathematical content from an unproved bridge. The procedures prescribe these construction responsibilities without specifying a complete automatic prose-generation algorithm. They apply to the explicitly activated Yang--Mills instance; they are not a default interpretation of every scientific paper. Their ordering is not inferred merely from appearing together in a binding list.

As a constructed paper example, a finite-block certificate checks a normalized bound for declared boundary data, while a volume-uniform supremum remains open. The narrative states the finite result; the proof appendix exposes the exact parameters and checked inequality; the open supremum stays a named obligation. An abstract claiming uniform control of all blocks would change the result even if the finite computation were reproduced perfectly. Faithful construction makes the useful certificate intelligible without pretending that the broader implication has been proved.

Original accepted status accepts the presentation; stopped status stops. Native reasoning only returns that disposition and does not publish a paper or execute an acceptance branch. The manifest leaves materialization blocked on the unresolved presentation projection. Thus the workflow describes a scientific communication boundary with an explicit native realization gap, rather than an automatic publication mechanism or a further proof step.

% SPDX-FileCopyrightText: © 2026 Idiap Research Institute <contact@idiap.ch>
%
% SPDX-FileContributor: Andre Freitas <andre.freitas@idiap.ch>
% Neuro-symbolic AI Group
%
% SPDX-FileContributor: Danilo Gusicuma <danilo.gusicuma@idiap.ch>
% Neuro-symbolic AI Group
%
% SPDX-License-Identifier: MIT
% Generated from step-explanations.json
\paragraph{Step-by-step reconstruction.}
An audit supports a finite-block certificate for declared boundary data but leaves a volume-uniform supremum estimate open.
The following is a source-guided illustrative reconstruction. Internal procedures are prescriptions; the native declarations do not establish their execution.
\begin{enumerate}[leftmargin=*,label=\arabic*.]
\item \textbf{Construct a faithful presentation} (\texttt{claim\_\allowbreak{}bounded\_\allowbreak{}presentation}).
\begin{description}[style=standard,labelindent=0pt,leftmargin=5.5em,labelwidth=5em,labelsep=0.5em,align=left,font=\normalfont\bfseries,itemsep=0.15em,parsep=0pt,topsep=0.3em]
\item[Purpose] Present the finite result within its supported scope.
\item[Input] Independent audit reference; claim ceiling
\item[Method] The native presentation producer is unresolved. Linked writers organize accepted results, expose proof dependencies and retain failed routes and open obligations while rendering the argument.
\item[Check] Each theorem must state the hypotheses it uses and where the finite certificate stops.
\item[Output] Intended presentation: finite certificate; unresolved uniform supremum
\item[Next] Release assessment consumes the presentation; string references do not enforce preservation automatically.
\item[Example] State exactly which boundary data the certificate checks, show the finite inequality, and list the supremum estimate as a separate open application burden.
\end{description}
\item \textbf{Assess release fidelity} (\texttt{release\_\allowbreak{}disposition}).
\begin{description}[style=standard,labelindent=0pt,leftmargin=5.5em,labelwidth=5em,labelsep=0.5em,align=left,font=\normalfont\bfseries,itemsep=0.15em,parsep=0pt,topsep=0.3em]
\item[Purpose] Assess manuscript fidelity to the audit.
\item[Input] Presentation; retained audit reference and claim ceiling
\item[Method] Compare headline, theorem and application claims against the accepted finite statement and remaining obligations.
\item[Check] Readability and compilation cannot strengthen the source claim.
\item[Output] Release disposition: supportable presentation or stopped overstatement
\item[Next] Original acceptance is a policy; native return neither publishes nor proves the result.
\item[Example] Release only the finite-result narrative with the uniformity gap visible; stop an abstract announcing all-block control from that same certificate.
\end{description}
\end{enumerate}
\paragraph{Changed-input contrast.}
If the abstract changes the finite claim into uniform control over all blocks, the presentation must fail its claim boundary even when every finite calculation and LaTeX build succeeds.


\subsection{Result elaboration: extending a manuscript within its evidence}
\label{app:granular-manuscript_result_elaboration}
\texttt{manuscript\_result\_elaboration} combines four occurrences. Initial calibration extracts the scope of a draft; abduction proposes an extension within it; refutation tests that extension; final calibration joins the same extension with its validity assessment. The three relations distinguish their licenses: \texttt{Constrains\allowbreak Manuscript\allowbreak Extension} carries the scope restriction, \texttt{Tests\allowbreak Manuscript\allowbreak Extension} exposes the proposal to criticism, and \texttt{Calibrates\allowbreak Manuscript\allowbreak Readiness} supports a readiness judgment from proposal plus verdict.
\[
 e_0:P\to\operatorname{Generated}[Q],\quad a:Q\to\operatorname{Generated}[H],
\]
\[
 r:H\to\operatorname{Generated}[V],\quad e_1:J\to\operatorname{Generated}[S].
\]
\begingroup\raggedright
Here $P$ denotes \texttt{Paper\allowbreak Draft}, $Q$ denotes \texttt{Manuscript\allowbreak Result\allowbreak Scope}, $H$ denotes \texttt{Manuscript\allowbreak Extension\allowbreak Hypothesis}, $V$ denotes \texttt{Extension\allowbreak Validity\allowbreak Assessment}, $J$ denotes \texttt{Manuscript\allowbreak Readiness\allowbreak Input}, and $S$ denotes \texttt{Manuscript\allowbreak Readiness\allowbreak Assessment}. For an input draft $p$, write $q=e_0(p)$, $z=a(q.\mathrm{value})$, $h=z.\mathrm{value}$, and $v=r(h)$ after successful generation. The final call to $e_1$ consumes $J$ with \texttt{extension\_hypothesis}$=h$ and \texttt{validity\_assessment}$=v.\mathrm{value}$. Final $S$ contains disposition, integrated draft, expansion validation, and readiness decision. These nominal records do not automatically enforce scientific scope or dependent identity.\par\endgroup

\paragraph{Source-prescribed conceptual work.}
The result-elaboration skills first recover the draft's smallest honest contribution and classify each expansion marker by its mathematical role. A missing coefficient, finite-object definition, transfer theorem, uniform estimate, and explanatory paragraph require different work. Open-module decomposition separates exact algebra or finite matrix checking from generative construction of estimates and proof strategies. The planner orders needs by genuine dependence: coefficients precede sign tests, and transfer certificates precede stronger application claims. Finite-object formalization specifies the carrier, indices, equivalence or support rule, quantitative measure, worked non-example, and the lemma that may consume the object. Proof-packet derivation then states each inference from its definitions and hypotheses, localizing a missing lemma or countermodel when it fails. Prior-art and disqualifier review test novelty and scope; narrative integration explains the result in the argument. Thus the four native occurrences compress an evidence-sensitive programme of mathematical elaboration, not a single generic request to expand prose.

In a constructed finite-result manuscript, an open module requests a sign certificate for a parameter-dependent matrix. Elaboration first defines the matrix, parameter range, and error terms; coefficient extraction then supplies the quantities needed for an exact sign test. A proof packet explains how that test enters the stated lemma. If an application additionally requires a geometric transfer estimate, finite sign success leaves that estimate open. The manuscript receives a checkable finite contribution and a precise boundary, rather than a paragraph asserting that the application now follows.

Original \texttt{proof\_required} routes to theorem delivery, \texttt{paper\_ready} accepts, and stalled reasoning stops. Native code returns the readiness record; comments describe a non-executable adapter to \texttt{Proof\allowbreak Queue\allowbreak Item} and a separate acceptance gate. There is no unresolved \texttt{\_} occurrence, but generation, integration, and a ready label alone do not prove an extension or execute publication.

% SPDX-FileCopyrightText: © 2026 Idiap Research Institute <contact@idiap.ch>
%
% SPDX-FileContributor: Andre Freitas <andre.freitas@idiap.ch>
% Neuro-symbolic AI Group
%
% SPDX-FileContributor: Danilo Gusicuma <danilo.gusicuma@idiap.ch>
% Neuro-symbolic AI Group
%
% SPDX-License-Identifier: MIT
% Generated from step-explanations.json
\paragraph{Step-by-step reconstruction.}
A finite-result draft requests a sign certificate for the matrix with diagonal entries 1 and off-diagonal entries t, for absolute t at most one half. A geometric application still needs a separate transfer estimate.
The following is a source-guided illustrative reconstruction. Internal procedures are prescriptions; the native declarations do not establish their execution.
\begin{enumerate}[leftmargin=*,label=\arabic*.]
\item \textbf{Recover the actual expansion need} (\texttt{e0}).
\begin{description}[style=standard,labelindent=0pt,leftmargin=5.5em,labelwidth=5em,labelsep=0.5em,align=left,font=\normalfont\bfseries,itemsep=0.15em,parsep=0pt,topsep=0.3em]
\item[Purpose] Distinguish mathematical gaps from prose gaps.
\item[Input] Draft; stated contribution; expansion markers
\item[Method] Classify each request as a definition, coefficient, finite check, transfer theorem or explanatory paragraph and recover the smallest honest scope.
\item[Check] Do not treat all expansion markers as requests to assert stronger results.
\item[Output] Expansion scope: finite matrix certificate; separate application obligation
\item[Next] Extension formation consumes this scope; unresolved geometric work remains distinct.
\item[Example] Specify the two-by-two symmetric matrix, the range absolute t at most one half, and the still-unproved geometric transfer estimate.
\end{description}
\item \textbf{Propose a bounded extension} (\texttt{a0}).
\begin{description}[style=standard,labelindent=0pt,leftmargin=5.5em,labelwidth=5em,labelsep=0.5em,align=left,font=\normalfont\bfseries,itemsep=0.15em,parsep=0pt,topsep=0.3em]
\item[Purpose] Construct the missing finite certificate.
\item[Input] Matrix scope; parameter range; missing sign certificate
\item[Method] Define the carrier and entries, order coefficient extraction before the sign test, and state the lemma that will consume it.
\item[Check] A finite extension must expose its quantitative range and avoid assuming the application.
\item[Output] Extension hypothesis: exact eigenvalue-based positivity certificate
\item[Next] Refutation tests this proposed extension and its parameter boundary.
\item[Example] For the matrix with diagonal 1 and off-diagonal t, propose eigenvalues 1+t and 1-t and hence a lower bound of one half on the stated range.
\end{description}
\item \textbf{Test the extension's validity} (\texttt{r0}).
\begin{description}[style=standard,labelindent=0pt,leftmargin=5.5em,labelwidth=5em,labelsep=0.5em,align=left,font=\normalfont\bfseries,itemsep=0.15em,parsep=0pt,topsep=0.3em]
\item[Purpose] Test the finite result and its application boundary.
\item[Input] Proposed matrix certificate; range; intended use
\item[Method] Check the algebra and boundary values, seek a non-example, and localize any absent transfer lemma.
\item[Check] Finite sign success does not discharge a geometric comparison hypothesis.
\item[Output] Validity assessment: scoped matrix calculation; open transfer debt
\item[Next] Readiness joins this assessment with the same extension hypothesis.
\item[Example] At t=1/2 the eigenvalues are 3/2 and 1/2; at t=2 they are 3 and -1. The range matters, and neither calculation supplies geometric transfer.
\end{description}
\item \textbf{Integrate with its evidence boundary} (\texttt{e1}).
\begin{description}[style=standard,labelindent=0pt,leftmargin=5.5em,labelwidth=5em,labelsep=0.5em,align=left,font=\normalfont\bfseries,itemsep=0.15em,parsep=0pt,topsep=0.3em]
\item[Purpose] Assess inclusion while preserving proof obligations.
\item[Input] Extension hypothesis + its validity assessment
\item[Method] Explain how the finite lemma fits the argument, retain the counterexample and separate paper readiness from unproved application claims.
\item[Check] An integrated paragraph cannot turn a missing bridge into an accepted theorem.
\item[Output] Finite certificate; expansion validation; readiness decision
\item[Next] Original proof-required routing needs a ProofQueueItem adapter; native return does not execute it.
\item[Example] Add the exact finite positivity result and its range to the draft, while sending the geometric transfer estimate to further proof work instead of claiming the application.
\end{description}
\end{enumerate}
\paragraph{Changed-input contrast.}
If t is allowed to exceed 1, the smaller eigenvalue becomes negative and the proposed positivity extension fails. If the finite range is retained, positivity still cannot establish the missing geometric transfer.


